\section{Exkurs: Tetens-Gleichung bzw. Magnus Gleichung}

Der Partialdruck eines Gases $e$ innerhalb eines Gasgemisches kann durch Umstellung der  \href{https://de.wikipedia.org/wiki/Thermische_Zustandsgleichung_idealer_Gase}{allgemeinen Gasgleichung} 
\begin{equation}
    p \cdot V = n \cdot R_{m} \cdot T  \rightarrow   p = \frac{n}{V} \cdot R_m \cdot T
\end{equation}
mit folgender Formel näherungsweise berechnet werden:
\begin{equation}
    e = \rho \cdot R_m \cdot T
\end{equation}
mit
\begin{itemize}
    \item  $p$ : Druck $\text{Pa} = \frac{\text{N}}{\text{m}^2} = \frac{J}{\text{m} \cdot \text{m}^2} = \frac{J}{\text{m}^3}$ 
    \item  $V$: Volumen in $\text{m}^3$ 
    \item  $n$: Stoffmenge in Mol 
    \item  $R_{m}$: Gaskonstande 8,31446261815324 $\frac{\mathrm{J}}{\mathrm{mol} \ \mathrm{K}}$ 
    \item  $T$: Termodynamische Temperatur in Kelvin 
    \item $\vartheta$: Termodynamische Temperatur in °C mit $T = 273{,}15 +  \vartheta$
    \item $\rho$ : Dichte in $\frac{\text{kg}}{\text{m}^3}$
\end{itemize} 
Für die Dampdichte $\rho_v$ folgt:
\begin{equation}
    \rho_v = \frac{p}{R_v \cdot T}
\end{equation}
Mit der Gaskonstante für Wasserdampf: $R = 461,52\,\frac{\text{J}}{\text{kg}\,\text{K}}$ und Sättigungsdampfdruck $\rho(T)$ folgt für den Druck  $p$
\begin{equation}
    p = \varphi \cdot \rho(T)
\end{equation}

Damit kann die \textbf{absolute Feuchte} der Luft in $\text{g} / \text{m}^3$ mit der folgenden Formel berechnet werden:
\begin{equation}
    \rho_v(\varphi, T) = \frac{\varphi \cdot p_s(T)}{R_v \cdot T_K}
\end{equation}
mit 
\begin{itemize}
    \item $\varphi$ : relative Luftfeuchtigkeit im Wertebereich $[0 \dots  1]$
    \item $p_s(T)$ :  Sättigungsdampfdruck in Pa (*saturation*)
    \item $R_v = 461,52\, \frac{\text{J}}{\text{kg} \cdot \text{K}}$: Gaskonstante für Wasserdampf (*vapor*)
    \item $T_K = T_{{}^\circ\text{C}} + 273,15\, \text{K}$ : Temperatur in Kelvin
\end{itemize}

Der Sättigungsdampfdruck $p_s(T)$ kann mit der \href{https://www.chemie.de/lexikon/S%C3%A4ttigungsdampfdruck.html} {\textbf{Tetens-Gleichung} bzw. \textbf{Magnus Gleichung}} approximiert werden:
\begin{equation}
    p_s(\vartheta) = 611{,}2\, \text{Pa} \cdot e^{\dfrac{17{,}62 \cdot \vartheta}{\vartheta + 243{,}12^\circ\text{C}}}
\end{equation}
wobei $\vartheta$ die Temperatur in °C angegeben ist.
Es gilt: 
\begin{itemize}
    \item $-30\,{}^\circ\mathrm{C} \leq 70\,{}^\circ\mathrm{C}$
    \item $p_s(\vartheta = 0\,{}^\circ\mathrm{C}) = 6{,}11213\;\mathrm{hPa} = 611{,}213\;\mathrm{Pa}$
\end{itemize}

Die Konstanten variieren in der Literatur, z.\ B.:
\begin{equation}
    p_s(\vartheta) 
    = 
    611{,}153\, \text{Pa} \cdot e^{\frac{22{,}4433 \cdot \vartheta}{\vartheta + 272{,}186^\circ\text{C}}}
\end{equation}

